% Figure: first-order RL circuit driven by a voltage step. Two curves on a
% shared normalised time axis: the inductor current rising toward its final
% value as 1 - exp(-t/tau), and the fraction of the final magnetic energy
% stored, which rises as (1 - exp(-t/tau))^2 and therefore lags the current.
% pgfplots.
\begin{figure}[htbp]
\centering
\begin{tikzpicture}
\begin{axis}[
  width=12cm, height=7.5cm,
  xlabel={time $t/\tau$}, ylabel={fraction of final value},
  xmin=0, xmax=6, ymin=0, ymax=1.1,
  axis lines=left,
  tick label style={font=\scriptsize},
  label style={font=\small},
  legend pos=south east,
  legend style={font=\scriptsize},
  clip=false,
]
% inductor current: 1 - exp(-t)
\addplot[engblue, very thick, domain=0:6, samples=140] {1 - exp(-x)};
\addlegendentry{current $i_L/I_\infty = 1 - e^{-t/\tau}$}
% stored energy fraction: (1 - exp(-t))^2
\addplot[engorange, very thick, dashed, domain=0:6, samples=140] {(1 - exp(-x))^2};
\addlegendentry{energy $U_L/U_\infty = (1 - e^{-t/\tau})^2$}
% asymptote
\addplot[enggray, dotted] coordinates {(0,1) (6,1)};
% one tau on current: 0.632
\draw[engred, dashed] (axis cs:1,0) -- (axis cs:1,0.632);
\filldraw[engblue] (axis cs:1,0.632) circle [radius=1.6pt];
\node[engblue, font=\scriptsize, anchor=north west] at (axis cs:1.05,0.62) {$i_L=63\%$ at $\tau$};
% one tau on energy: 0.400
\filldraw[engorange] (axis cs:1,0.400) circle [radius=1.6pt];
\node[engorange, font=\scriptsize, anchor=north west] at (axis cs:1.05,0.38) {$U_L=40\%$ at $\tau$};
\end{axis}
\end{tikzpicture}
\caption[RL current and stored energy]{Step response of a first-order RL
circuit. The inductor current rises toward its final value with time constant
$\tau = L/R$, reaching $63\%$ after one time constant. The stored magnetic
energy $U_L = \tfrac{1}{2}Li_L^2$ tracks the square of the current, so it lags
the current and stands at only $40\%$ of its final value when the current has
reached $63\%$. Sizing an inductor for a settling time therefore leaves it well
short of its final stored energy at one time constant, the reason a switching
converter runs several time constants of margin before it counts an inductor
as charged.}
\label{fig:vol04ch09:rl-energy}
\end{figure}
